% Reference draft only: rewrite this paper in the student's own voice.
% This is a separate DSCI602 Option C reference draft, not the root paper.
\documentclass[10pt,conference]{IEEEtran}

\usepackage{amsmath}
\usepackage{booktabs}
\usepackage{cite}
\usepackage[hidelinks]{hyperref}
\usepackage{microtype}

\title{Fairness-Aware Service Allocation in Adaptive Quantum Networks:\\
Risks, Metrics, and Auditable Mediation}

\author{\IEEEauthorblockN{Piter Garcia}
\IEEEauthorblockA{DSCI 602: Applied Data Science II\\
Focus Area: Fairness\\
Capstone: Fairness-Aware Adaptive Routing for Quantum Networks}}

\begin{document}
\maketitle

\begin{abstract}
This reference analysis examines fairness in a quantum-network controller that
selects entanglement routes and allocates scarce qubits under changing threat
conditions. The central concern is allocative: an efficiency-oriented bandit
policy may deliver reliable service to some flow classes while concentrating
delay, retries, or starvation on others. The proposed response is not a claim
that one fairness statistic solves the problem. It is a measurement-first
workflow that defines service classes before evaluation, records demand and
resource access, compares an unmediated policy with an auditable mediator, and
reports the resulting utility--disparity frontier with explicit residual risk.
\end{abstract}

\section{System Context and Responsibility Boundary}
The artifact under analysis is a threat-aware quantum-network controller, not a
generic classifier. At each decision round, the controller observes a network
state, selects an entanglement route, and allocates limited qubit-generation
capacity. Link quality, memory constraints, replay capacity, allocator choice,
and adversarial conditions can change the quality of the available actions.
Prior routing work supplies the inherited bandit and threat-regime evaluation
spine; the DSCI 602 contribution is a planned fairness-aware mediation layer
for quantum service equity. Routing and resource allocation therefore remain
the primary technical problem, while clinical applications are future context
only.

The system influences who receives successful and timely quantum service under
scarcity. Relevant stakeholders are applications or tenants competing for
flows, network operators responsible for allocation policy, and downstream
users who depend on the resulting entanglement paths. The initial equity units
should be operational service classes defined before inspecting outcomes, such
as demand tiers, latency requirements, or context-quality cohorts. These are
not assumed to be demographic groups; if a future deployment maps them to
people or institutions, protected-attribute analysis and stakeholder review
would be required.

This boundary matters because a network-wide efficiency score can conceal
service differences. Inherited configuration-specific evidence includes a
Thompson-sampling replay slice whose Oracle-normalized efficiency changes from
89.2\% to 84.9\% and then 90.8\% as replay capacity increases under an
online-adaptive threat. That sequence is not a DSCI 602 fairness result. It is
an example of why aggregate behavior deserves disaggregation: a gain or loss
may be distributed unevenly across flow classes. Quantum-network routing and
fair-sharing research likewise treats scarce capacity and simultaneous
requests as allocation questions rather than merely throughput questions
\cite{li2021,cicconetti2023}.

\section{Where Fairness Risk Enters}
Suresh and Guttag's life-cycle framework provides a useful map for this
pipeline \cite{suresh2021}. Representation harm can occur when evaluation
omits difficult topologies, low-context flows, or relevant service classes.
Measurement harm occurs when traces record success and reward but omit demand,
allocated qubits, route choice, fidelity, retries, starvation age, or context
quality. Aggregation harm occurs when one reward treats incompatible service
needs as interchangeable. Evaluation harm occurs when aggregate efficiency is
reported without class-level outcomes. Deployment harm occurs when a policy
tuned at one threat regime is transferred to another operating point.

Sequential learning adds a feedback-loop risk. A bandit observes evidence only
for selected actions. If a route or class receives fewer opportunities, it can
produce less reward evidence, which can reduce its future selection probability.
Fairness constraints therefore have a learning and regret cost, rather than
being a free post-processing step \cite{joseph2016}. Fair contextual bandit
work similarly motivates minimum service rates and explicit treatment of
context when outcomes differ across users \cite{chen2020}. In this system,
poor context can become an allocation disadvantage even when no protected
attribute is used.

The plausible harms are concrete. High-demand or high-degree flows may receive
most successful allocations while low-demand flows absorb waiting and retries.
Flows with weak context may be treated as unreliable and consequently receive
less exploration. Long or low-fidelity routes may be repeatedly deprioritized,
creating starvation that disappears inside a network-wide mean. These are
allocative harms: access to a scarce capability is withheld or degraded. They
also create accountability harms if operators cannot reconstruct why a flow
received inferior service. The general fairness literature shows why aggregate
performance is insufficient: class-level failures can be hidden by overall
accuracy or a seemingly neutral proxy \cite{buolamwini2018,obermeyer2019}.

\section{Fairness Objective and Measurement}
Classical definitions are useful warnings but cannot be copied mechanically
into a resource-allocation system. Demographic parity would compare success or
access rates across classes, while equalized odds and equality of opportunity
compare error behavior conditional on an outcome \cite{hardt2016}. Calibration
and parity can conflict under unequal base rates, and impossibility results
show that several desirable criteria cannot generally hold simultaneously
\cite{chouldechova2017,kleinberg2017}. The responsible objective is therefore
to report a frontier rather than declare one scalar to be ``fair.''

One candidate primary measure is demand-normalized access. For class $g$ over
a run horizon $T$, let $Q_{g,t}$ be allocated quantum capacity and $D_{g,t}$
be declared demand:
\begin{equation}
 A_g(T)=\frac{\sum_{t=1}^{T}Q_{g,t}}
 {\sum_{t=1}^{T}D_{g,t}}, \qquad
 \Delta_A(T)=\max_{g,h}|A_g(T)-A_h(T)|.
\end{equation}
The audit should also report success-rate and latency gaps, worst-class
outcomes, retry and starvation frequency, fidelity failures, and a bounded
resource-allocation index such as Jain's index \cite{jain1984}. Zero-demand
classes require a declared exclusion rule, and every class definition,
denominator, physical-infeasibility condition, and missing-context rule should
be versioned with the experiment. The point is not to hide complexity inside a
single index; it is to make the trade-offs inspectable.

\section{Mediation and Evaluation Protocol}
The first mitigation is measurement. An enriched decision trace should record
the class definition, demand, context-quality indicators, candidate routes,
policy scores, selected route, allocated qubits, fidelity estimate, success,
latency, retries, and starvation age. Constructed unit tests should include
equal service, unequal service, zero demand, missing context, and physically
infeasible routes. This is a prerequisite for any claim about disparity.

The second mitigation is an in-processing mediator attached at a documented
policy-score or allocator boundary. A candidate score can be written as
\begin{equation}
 S_t^{\mathrm{fair}}(a)=U_t(a)-\lambda[\widehat{\Delta}_t(a)-\epsilon]_+
 -\gamma C_t(a),
\end{equation}
where $U_t$ is base utility, $\widehat{\Delta}_t$ is estimated disparity,
$\epsilon$ is tolerated slack, and $C_t$ penalizes missing or unreliable
context. The identity setting $\lambda=\gamma=0$ is an explicit off-switch,
not an afterthought. The mediator should remain separable and auditable so
that the inherited routing policy can be compared with its fairness-aware
counterpart \cite{chen2020}.

Validation must compare the unmediated baseline with a prespecified sweep of
$\lambda$, $\gamma$, and $\epsilon$ across matched policies, allocators,
capacities, topologies, threat regimes, and random seeds. Report aggregate
utility, regret, success, and fidelity beside $\Delta_A$, latency gaps,
worst-class outcomes, Jain's index, and starvation frequency. Repeated trials
and uncertainty intervals are necessary because small classes can produce
unstable estimates. The result should be a utility--disparity frontier, not a
single favorable operating point selected after inspection. If a class has too
few observations for a stable estimate, the system should report uncertainty
and apply a conservative service floor rather than infer that no disparity
exists.

\section{Limitations and Residual Risk}
This analysis is prospective. It does not claim completed live mediation or
completed DSCI 602 fairness experiments. The simulator does not establish
hardware-grounded timing, contention, or noise behavior. Service classes are
design choices and may encode the wrong units of concern; demand itself may be
strategically inflated. A mediator can reduce utility, increase latency for
other classes, or favor an imperfect proxy for equity. Conversely, removing
mediation can preserve aggregate utility while hiding persistent exclusion.

The remaining governance questions are material: who approves classes and
denominators, who reviews alerts, who can change fairness weights, and what
rollback path exists? Security and privacy intersections, including poisoning
of threat feedback and exposure of flow telemetry, require separate analysis.
The responsible claim is narrow: adaptive quantum routing can create
service inequity while appearing efficient, so the project should measure
flow-conditioned outcomes, evaluate an auditable mediator, and disclose the
frontier and its residual uncertainty.

\bibliographystyle{IEEEtran}
\bibliography{refs}

\end{document}
